\section{Introducción: motivación para la guía en tiempo de inferencia}

En las semanas anteriores del curso, discutimos los procesos básicos de entrenamiento y generación de modelos de difusión (Diffusion Model) y modelos de flujo (Flow Model), así como cómo acelerar la generación mediante técnicas en tiempo de inferencia o ajuste fino. El problema central de esta clase gira en torno a otra dirección: \textbf{¿cómo aprovechar los modelos preentrenados para inyectar señales de guía durante la inferencia, de modo que los resultados generados cumplan preferencias o condiciones específicas del usuario?}

\begin{importantbox}{Problema central de la lección}
Dado un modelo de difusión/flujo preentrenado (sin condiciones), ¿cómo podemos guiar el proceso de generación en tiempo de inferencia sin volver a entrenar, para que los puntos de datos resultantes satisfagan condiciones específicas —por ejemplo, coincidan con observaciones parciales dadas, sigan ciertas restricciones de disposición o de estilo?
\end{importantbox}

Actualmente existen numerosos modelos de difusión preentrenados que abarcan múltiples modalidades, como generación de imágenes, generación de video, diseño molecular y planificación del movimiento robótico. Estos modelos suelen requerir grandes volúmenes de datos y capacidad computacional para su entrenamiento, por lo que es difícil para investigadores comunes entrenarlos desde cero. Por ello, surge una pregunta natural: \textbf{¿es posible hacer más cosas utilizando únicamente los modelos preentrenados?}

Una ventaja clave de los modelos de difusión/flujo radica en su característica de \textbf{generación iterativa}. A diferencia de los modelos de generación de un solo paso como GAN y VAE, en los modelos de difusión, a lo largo de la trayectoria desde el ruido $x_T$ hasta los datos limpios $x_0$, cada paso puede ser "manipulado": podemos modificar la velocidad predicha o el ruido para curvar la trayectoria hacia el punto final deseado.

\begin{knowledgebox}{Repaso: Classifier Guidance y Classifier-Free Guidance}
En clases anteriores, ya hemos visto dos métodos para inyectar guía durante el proceso de generación:
\begin{itemize}
    \item \textbf{Classifier Guidance}: Entrena una red neuronal adicional de clasificación que utiliza sus gradientes sobre los datos de ruido intermedios $x_t$ para proporcionar un puntaje de verosimilitud (likelihood score), ajustando así la trayectoria de generación.
    \item \textbf{Classifier-Free Guidance (CFG)}: Utiliza la misma red para aprender simultáneamente el ruido condicional y el incondicional, calculando implícitamente el puntaje de verosimilitud a través de su diferencia y controlando la intensidad de la señal de guía mediante un parámetro $w$.
\end{itemize}
Ambos métodos requieren \textbf{entrenamiento}: ya sea entrenar un clasificador adicional o incluir pares de datos condicionales e incondicionales durante el entrenamiento del predictor de ruido. Esta lección aborda métodos de guía en tiempo de inferencia que \textbf{no requieren entrenamiento adicional}.
\end{knowledgebox}

\subsection{Tres categorías principales de guía en tiempo de inferencia}

El ponente clasifica las técnicas de guía en tiempo de inferencia en tres grandes grupos:

\begin{enumerate}
    \item \textbf{Score Manipulation (manipulación de puntajes)}: Modificaciones prácticas directas a la función de puntajes durante el proceso de desruido, sin garantías teóricas estrictas, pero con buenos resultados en la práctica.
    \item \textbf{Guía basada en marcos de problemas inversos}: Modela la generación condicional como un problema inverso (Inverse Problem), utilizando la fórmula de Bayes para aproximar el cálculo del puntaje de verosimilitud, lo que otorga cierta base teórica.
    \item \textbf{Otras técnicas avanzadas}: Se discutirán en la próxima lección (Lecture 12).
\end{enumerate}

Esta lección cubre principalmente las primeras dos categorías de métodos.

\subsection{Resumen de la sección}

La generación guiada en tiempo de inferencia es una ventaja central de los modelos de difusión, derivada de su naturaleza de generación iterativa. En esta lección, partiremos desde técnicas prácticas de manipulación de puntajes para adentrarnos progresivamente en marcos de resolución de problemas inversos con soporte teórico.



